\chapter{Integration et deploiement}

\section{Integration inter-projets}

\subsection*{Frontend <-> Backend}
Le frontend Flutter consomme l'API backend sur:
\begin{itemize}[leftmargin=*]
  \item REST: \texttt{http://<host>:8091/football-academy/api}
  \item WebSocket STOMP: \texttt{ws://<host>:8091/football-academy/ws}
\end{itemize}

\subsection*{Frontend <-> Chatbot}
Le module chatbot Flutter appelle le service Django sur \texttt{http://localhost:8000/api/chat} (ou \texttt{10.0.2.2} sur emulateur Android).

\section{Variables de configuration}

\begin{table}[H]
\centering
\begin{tabular}{p{0.35\textwidth} p{0.55\textwidth}}
\toprule
\textbf{Variable} & \textbf{Usage} \\
\midrule
\texttt{PORT} & Port backend Spring Boot (defaut 8091). \\
\texttt{DB\_URL, DB\_USER, DB\_PASS} & Connexion base de donnees backend. \\
\texttt{CHATBOT\_API\_KEY} & Protection optionnelle endpoint chatbot. \\
\texttt{QA\_CSV} & Dataset Q/R du chatbot. \\
\texttt{MIN\_SIM, FUZZY\_MIN} & Parametrage du moteur de similarite chatbot. \\
\bottomrule
\end{tabular}
\caption{Variables principales de deploiement}
\end{table}

\section{Strategie de deploiement recommandee}

\subsection*{Backend}
Deploiement conteneurise via Docker, pilote par pipeline CI/CD (test/build/docker push).

\subsection*{Frontend mobile}
Build et distribution independants pour mobile (Android/iOS).

\subsection*{Frontend admin web}
Build web separe et hebergement sur une plateforme statique ou serveur web. Cette separation evite de lier les contraintes web admin au package mobile.

\subsection*{Chatbot}
Deploiement en service Python separable (VM/container) selon la charge. En local, execution possible sur la meme machine que le frontend.

\section{Procedure de lancement local}

\subsection*{Backend Spring Boot}
\begin{lstlisting}[style=codeblock,language=bash]
cd football-academy
mvn clean package
mvn spring-boot:run
\end{lstlisting}

\subsection*{Frontend Flutter}
\begin{lstlisting}[style=codeblock,language=bash]
cd moez_project
flutter pub get
flutter run
\end{lstlisting}

\subsection*{Chatbot Django}
\begin{lstlisting}[style=codeblock,language=bash]
cd chatbot/chatbot
python -m venv env_py10
# activation env selon systeme
pip install -r requirements.txt
python manage.py migrate
python manage.py runserver 127.0.0.1:8000
\end{lstlisting}

\section{Points de vigilance en production}

\begin{itemize}[leftmargin=*]
  \item externaliser les secrets (JWT, DB, API keys) hors du code source;
  \item verifier les CORS et restrictions de domaines autorises;
  \item monitorer les performances chat temps reel et chatbot;
  \item mettre en place des sauvegardes DB et retention logs.
\end{itemize}
